\section{Conclusion}

The original 288-cell audit established exact representability but confounded
the generating equation, graph information, and head form.  The retrospective
diagnostic exposed that defect; it did not solve it.

Two newly frozen prospective cohorts do.  In the attribution cohort, train-only
selection recovered two-edge graphs and head families under stochastic noise.
Correct graph information improved held-out two-mechanism composition NLL under
factorized, parameter-matched generic sparse, and larger dense heads.  The
seven-sparse embedding proposition establishes an equivalent generic support,
and the sparse procedure recovered that support in every world, producing an
exact tie with the factorized head.  The distinctive object is therefore the
typed minimal support and its learned graph, not surface notation.

The population result now identifies that motif and its typed heads under
full-support primitive interventions and symmetric corruption below \(6/7\);
the finite bound states the repeated-input coverage required for modal recovery.
The random-design greedy selector is not covered by that theorem, but its
120/120 recovery is stable across the post-review \(3\times3\) noise grid.

In the separate stress cohort, an unmodeled nonadditive synergy made the fitted
model nonexact in every world, yet the learned graph retained an NLL effect of
0.18545 over wrong routing.  This removes exact generator--predictor equality as
an explanation of the graph effect.  A clean-room second-language
implementation reproduced every learned fit and analysis mean in the main
cohort.  Within the stated finite stochastic population, graph information,
minimal structural support, composition, misspecification, and post-export
fitting/evaluation/analysis implementation reproducibility are therefore
separately identified.

The next decisive comparison is not a factorized head versus a generic sparse
head. It contrasts two models that have both recovered the same dependency
graph under entity creation and deletion, identity maintenance, and relation
reconfiguration. For that comparison to identify anything, both arms must
receive identical observable information and differ only in representational
commitment. Supplying identity or typed relations as input to one arm only
would make the result an artifact of information asymmetry, the same class of
confound diagnosed in the original finite-family audit. The turnover quantity
\(b(p)\), its zero criterion, and the vacuity proposition for survivor-only
preservation are already formalized in the companion theory; a turnover-focused
study is the smallest self-contained instance of this design.
